\chapter{Control flow}
\label{chap:control}

Every program of the previous chapters did the same thing on every run: it
started at the top of \token{main}, went through each instruction once, in the
order written, and stopped at the bottom. A recipe that never says ``if the
dough is too dry, add water'' or ``stir until smooth'' can only describe the
simplest dishes, and the same is true of a program.

This chapter gives the program those two abilities. It can \emph{decide}, taking
one path or another depending on a value, and it can \emph{repeat}, running the
same instructions again as long as there is work left. Together they are called
\textit{control flow}, because they choose where control -- the cursor of
Section~\ref{sec:types:variable_declaration} -- goes next. By the end of the
chapter, a few dozen lines are enough to write a small game.

Part II covers the same constructions exhaustively, in
Chapter~\ref{chap:flow}.

\minitoc%

\vfill\pagebreak
\section{Values known only at run time}
\label{sec:control:runtime_values}

A program decides by asking a question about a value, such as ``is the age at
least 18?'', and the answer chooses what it does next. In the programs so far,
though, every value was written in the source code. If a program contains
\token{let age = 41;}, the question ``is the age at least 18?'' gets the same
answer every time the program runs: yes. Asking it is pointless, since the
program could print its answer directly. Ymir goes further and refuses such a
question outright, as Section~\ref{sec:control:if} shows.

A decision is only useful when its answer can change from one run to the next.
For that, the program needs values that are not written in its source code,
values that it learns only while it runs.

The simplest source of such values is the person running the program. The
function \token{read}, from the module \token{std::io} that already provides
\token{println}, waits for the user to type a line and returns what was typed.

\begin{lstlisting}[style=coloredverbatim, caption=Reading a number typed by the user, label=lst:control:read_age]
use std::io;

fn main() {
  let age = read!i32(ask-> "How old are you? ");
  println("Next year you will be ", age + 1, ".");
}
\end{lstlisting}
\vspace{-10pt}%
\begin{lstlisting}[style=bashVerb, escapechar=@+]
@+\textcolor{teal!80}{alice@dev:\mtilde}@+$ gyc main.yr -o main
@+\textcolor{teal!80}{alice@dev:\mtilde}@+$ ./main
How old are you? 41
Next year you will be 42.
\end{lstlisting}

The call on line~4 has two parts that the book has not presented yet. Until
Section~\ref{sec:functions:read} explains them, treat the whole expression as a
fixed formula:

\begin{itemize}
\item \token{!i32} says which type of value to read. Any other integer type
  works in its place; this chapter reads only \token{i32}.
\item \token{ask-> "..."} gives the text to print before waiting, so the user
  knows what is expected. It can be left out, as in \token{read!i32()}.
\end{itemize}

\token{read} does not check what the user types. If the line does not start
with a number, \token{read!i32} returns \token{0}, and the program carries on
with it. The
programs of this chapter are small enough to live with that.

The second source of values is chance. The function \token{uniform}, from the
module \token{std::rand}, returns a number picked at random between two bounds,
both included. The program below prints a different number, from 1 to 6, each
time it runs.

\begin{lstlisting}[style=coloredverbatim, caption=Rolling a die, label=lst:control:die]
use std::io;
use std::rand;

fn main() {
  let die = uniform(1, 6);
  println("The die shows ", die, ".");
}
\end{lstlisting}

With \token{read} and \token{uniform}, a program can hold values that the
compiler cannot know in advance. The rest of the chapter uses them to make
decisions.

\vfill\pagebreak
\section{Blocks}
\label{sec:control:blocks}

Every construction of this chapter works on \textit{blocks}, so they come first.
A block is a sequence of instructions between braces, \token{\{} and \token{\}}.
The body of \token{main} is one. A block can also stand anywhere a value is
expected, and then it does two things.

\begin{itemize}
\item It \textbf{groups} instructions. They run in order, from the opening
  brace to the closing one, and the variables declared inside exist only until
  the closing brace. Past that point their names mean nothing.
\item It \textbf{has a value}: the value of its last expression, provided that
  expression is not followed by a semicolon.
\end{itemize}

In the next listing, the block on lines~5 to~8 computes a price with a tax of
one fifth. The variable \token{tax} is a step of the calculation and nothing
more, so it lives inside the block. The block's value is \token{price + tax},
written last and without a semicolon, and that is what \token{total} receives.

\begin{lstlisting}[style=coloredverbatim, caption=A block used as a value, label=lst:control:block_value]
use std::io;

fn main() {
  let price = read!i32(ask-> "Price: ");
  let total = {
    let tax = price / 5;
    price + tax
  };
  println("Total with tax: ", total);
}
\end{lstlisting}
\vspace{-10pt}%
\begin{lstlisting}[style=bashVerb, escapechar=@+]
@+\textcolor{teal!80}{alice@dev:\mtilde}@+$ ./main
Price: 100
Total with tax: 120
\end{lstlisting}

The two properties each have a way to go wrong. Once the block is closed,
\token{tax} no longer exists, and a line~9 that tried to print it would stop the
compilation with ``undefined symbol tax''.

The other mistake is to end the last expression with a semicolon. The semicolon
turns an expression into an instruction whose value is thrown away, so the block
no longer has a value to give. Its type is then \token{void}, the type of
``nothing'', and no variable can hold nothing.

\begin{lstlisting}[style=coloredverbatimError, escapechar=@, caption=A block whose last expression ends with a semicolon]
use std::io;

fn main() {
  let price = read!i32(ask-> "Price: ");
  let @\hb{total}@ = {
    let tax = price / 5;
    price + tax;
  };
  println("Total with tax: ", total);
}
\end{lstlisting}

The compiler reports that it ``cannot declare a variable of type void''. The
line it points at is the declaration of \token{total}, but the fault is the
semicolon on line~7.

A block does not have to give a value. It can also stand where an instruction is
expected, between two other instructions, and then it only groups. That is
useful on its own: a variable needed for a few lines, declared in such a block,
disappears at the closing brace, and the rest of the function cannot use it by
mistake. Blocks are used most, though, in the constructions that follow, where
each branch of a decision and each repetition of a loop is a block.

\vfill\pagebreak
\section{Making a decision: \texttt{if} and \texttt{else}}
\label{sec:control:if}

\token{if} runs a block only when a condition holds. The condition is any
expression of type \token{bool}, the type of Section~\ref{sec:types:boolean},
and the block runs when it is \token{true}. When it is \token{false}, the block
is skipped, and the program continues after it.

\begin{lstlisting}[style=coloredverbatim, caption=A block run only when a condition holds, label=lst:control:if]
use std::io;

fn main() {
  let temperature = read!i32(ask-> "Temperature: ");
  if temperature > 30 {
    println("It is hot, drink some water.");
  }
  println("Have a nice day.");
}
\end{lstlisting}
\vspace{-10pt}%
\begin{lstlisting}[style=bashVerb, escapechar=@+]
@+\textcolor{teal!80}{alice@dev:\mtilde}@+$ ./main
Temperature: 35
It is hot, drink some water.
Have a nice day.
@+\textcolor{teal!80}{alice@dev:\mtilde}@+$ ./main
Temperature: 12
Have a nice day.
\end{lstlisting}

The last line is printed on both runs, because it is outside the block: the
\token{if} decides about line~6 only.

\subsection{Two paths: \texttt{else}}

\token{else}, written after the block of an \token{if}, adds a second block that
runs when the condition is \token{false}. One of the two blocks always runs, and
never both. The next listing tells even numbers from odd ones with the remainder
operator \tokennolst{\%}, and Figure~\ref{fig:control:if_else} draws the two paths it
can take.

\begin{lstlisting}[style=coloredverbatim, caption=Two branches, label=lst:control:if_else]
use std::io;

fn main() {
  let n = read!i32(ask-> "A number: ");
  if n % 2 == 0 {
    println(n, " is even.");
  } else {
    println(n, " is odd.");
  }
}
\end{lstlisting}

\begin{figure}[h]
  \centering
  \begin{adjustbox}{max width=\linewidth}
    \begin{tikzpicture}[node distance=6mm and 4mm]

      \node[cfterm] (start) {start of main};
      \node[cfstep, below=of start] (read) {let n = read!i32(...);};
      \node[cftest, below=of read] (test) {n \% 2 == 0};
      \node[cfstep] (then) at ($(test) + (-32mm, -16mm)$) {println(n, " is even.");};
      \node[cfstep] (else) at ($(test) + (32mm, -16mm)$) {println(n, " is odd.");};
      \node[cfterm] (end) at ($(test) + (0, -30mm)$) {end of main};

      \draw[cfflow] (start) -- (read);
      \draw[cfflow] (read) -- (test);
      \draw[cfflow] (test.west) -| node[cflab, above, pos=0.4] {true} (then.north);
      \draw[cfflow] (test.east) -| node[cflab, above, pos=0.4] {false} (else.north);
      \draw[cfflow] (then.south) |- (end.west);
      \draw[cfflow] (else.south) |- (end.east);
    \end{tikzpicture}
  \end{adjustbox}
  \caption{\label{fig:control:if_else} The two paths through
    Listing~\ref{lst:control:if_else}. Exactly one of the two branches runs,
    then both paths meet again.}
\end{figure}


\subsection{More than two paths: \texttt{else if}}

The block after \token{else} can be replaced by another \token{if}, which gives a
chain of conditions. They are tested from top to bottom, and the first one that
holds wins: its block runs, and the rest of the chain is skipped. The final
\token{else} catches every case that no condition matched.

\begin{lstlisting}[style=coloredverbatim, caption=A chain of conditions, label=lst:control:else_if]
use std::io;

fn main() {
  let score = read!i32(ask-> "Score (0 to 100): ");
  if score >= 90 {
    println("Grade A");
  } else if score >= 75 {
    println("Grade B");
  } else if score >= 50 {
    println("Grade C");
  } else {
    println("Grade F");
  }
}
\end{lstlisting}

The order does the work here. A score of 95 is also above 75 and above 50, but
the first test catches it, so it gets an A and nothing else. Written in the
opposite order, starting with \token{score >= 50}, the chain would give a C to
everyone who passes.

The final \token{else} is optional, as it is after a single \token{if}. Without
it, a value that meets none of the conditions runs none of the blocks, and the
program simply continues after the chain. The next listing warns about the two
extremes of temperature and says nothing about the others.

\begin{lstlisting}[style=coloredverbatim, caption=A chain with no final \token{else}, label=lst:control:else_if_no_else]
use std::io;

fn main() {
  let temperature = read!i32(ask-> "Temperature: ");
  if temperature > 30 {
    println("It is hot, drink some water.");
  } else if temperature < 0 {
    println("It is freezing, wear a coat.");
  }
  println("Have a nice day.");
}
\end{lstlisting}
\vspace{-10pt}%
\begin{lstlisting}[style=bashVerb, escapechar=@+]
@+\textcolor{teal!80}{alice@dev:\mtilde}@+$ ./main
Temperature: -5
It is freezing, wear a coat.
Have a nice day.
@+\textcolor{teal!80}{alice@dev:\mtilde}@+$ ./main
Temperature: 15
Have a nice day.
\end{lstlisting}

\subsection{Choosing a value}

An \token{if} with an \token{else} is also an expression, whose value is the
value of the block that ran. Section~\ref{sec:control:blocks} gave each block a
value, so a decision between two values can be written in one line.

\begin{lstlisting}[style=coloredverbatim, caption=An \token{if} used as a value, label=lst:control:if_value]
use std::io;

fn main() {
  let n = read!i32(ask-> "A number: ");
  let parity = if n % 2 == 0 { "even" } else { "odd" };
  println(n, " is ", parity, ".");
}
\end{lstlisting}

For this to make sense, both blocks must produce a value, of the same type.
Without an \token{else}, the program would have no value to give
\token{parity} when \token{n} is odd, and the compiler refuses it:
\errmsg{E4081}{the if condition has no else value, so it must produce a void
  value not a mut [c8]}.

\begin{lstlisting}[style=coloredverbatimError, escapechar=@, caption=An \token{if} with no \token{else} has no value]
use std::io;

fn main() {
  let n = read!i32(ask-> "A number: ");
  let parity = @\hb{if n \% 2 == 0 \{ "even" \}}@;
  println(n, " is ", parity, ".");
}
\end{lstlisting}

The same rule rejects \token{if n > 0 \{ 1 \} else \{ "no" \}}, whose two
blocks produce an \token{i32} and a string: the type of a variable is fixed when
it is declared (Section~\ref{sec:types:variable_declaration}), so it cannot hold
one or the other depending on the run.

\subsection{Conditions the compiler refuses}

A condition must be a \token{bool}. Some languages accept a number and treat
zero as false; Ymir does not, and \token{if n \{} is an error,
\errmsg{E4095}{incompatible types i32 and bool}. The test has to be spelled out: \token{if n != 0 \{}.

A condition must also be unknown at compile time. In the next listing,
\token{limit} is a constant equal to 30, so the compiler knows the test is
always \token{true}: the \token{if} decides nothing, and the compiler says so
rather than let it pass.

\begin{lstlisting}[style=coloredverbatimError, escapechar=@, caption=A decision the compiler can already make]
use std::io;

fn main() {
  let limit = 30;
  if @\hb{limit > 20}@ {
    println("The limit is above 20.");
  }
}
\end{lstlisting}

The message is \errmsg{E4268}{conditional test is always evaluated to true,
  branch is always entered}. A condition that is always false gets the matching
message, E4266, since its
block could never run. A condition therefore has to depend on a value that the
compiler cannot know. In this section, that is always a value typed by the user
or drawn at random, with \token{read} and \token{uniform}, from
Section~\ref{sec:control:runtime_values}. The loops of the next sections add a
third kind: a variable that the program itself changes as it runs.

\notebox{The block of an \token{if} or of an \token{else} that holds a single
  instruction can be written without its braces: \token{if n < 0
  println("negative");} is valid. Not every block allows it. Apart from these
  two, only the \token{for} of Section~\ref{sec:control:for} and the arms of a
  \token{match} do; the body of a function, of a \token{loop} or of a
  \token{while} always needs its braces. The book writes them everywhere,
  because adding a second instruction to an \token{if} without braces silently
  puts it outside the \token{if}.}

\vfill\pagebreak
\section{Repeating: \texttt{loop}, \texttt{break} and \texttt{continue}}
\label{sec:control:loop}

\token{loop} runs a block over and over. When control reaches the closing brace,
it goes back to the opening one, and the block runs again. The \token{loop}
itself never stops, so the block must contain the instruction that ends it:
\token{break}, which leaves the loop at once. The program then continues with
the instruction that follows the loop.

The next listing adds up the numbers typed by the user, as many as there are. It
cannot know in advance how many that is, so it loops, and it breaks out when the
user types \token{0}.

\begin{lstlisting}[style=coloredverbatim, caption=Repeating until the user says stop, label=lst:control:loop]
use std::io;

fn main() {
  let mut total = 0;
  loop {
    let n = read!i32(ask-> "A number (0 to stop): ");
    if n == 0 {
      break;
    }
    total += n;
  }
  println("Total: ", total);
}
\end{lstlisting}
\vspace{-10pt}%
\begin{lstlisting}[style=bashVerb, escapechar=@+]
@+\textcolor{teal!80}{alice@dev:\mtilde}@+$ ./main
A number (0 to stop): 4
A number (0 to stop): 10
A number (0 to stop): 3
A number (0 to stop): 0
Total: 17
\end{lstlisting}

Each time round, the block declares a new \token{n}, which lives until the end of
that repetition, called an \textit{iteration}. \token{total} is declared before
the loop, so it survives from one iteration to the next and accumulates the sum.
Declared inside the block, it would start again at \token{0} on each iteration.

\tipbox{A loop whose \token{break} is never reached runs forever. If a program
  seems stuck, \textit{Ctrl+C} in the terminal stops it.}

\subsection{Skipping the rest of an iteration: \texttt{continue}}

\token{continue} is the other way to jump. Instead of leaving the loop, it
abandons the current iteration and goes straight back to the top of the block
for the next one. The next listing refuses negative numbers: for those, line~12
skips the addition on line~14. It also stops by itself once the total reaches
100, with a second \token{break} on line~16. A loop can have any number of
\token{break}s, and the first one reached ends it.
Figure~\ref{fig:control:loop_jumps} draws the three jumps.

\begin{lstlisting}[style=coloredverbatim, caption=Skipping an iteration, label=lst:control:continue]
use std::io;

fn main() {
  let mut total = 0;
  loop {
    let n = read!i32(ask-> "A number (0 to stop): ");
    if n == 0 {
      break;
    }
    if n < 0 {
      println("Negative numbers are ignored.");
      continue;
    }
    total += n;
    if total >= 100 {
      break;
    }
  }
  println("Total: ", total);
}
\end{lstlisting}
\vspace{-10pt}%
\begin{lstlisting}[style=bashVerb, escapechar=@+]
@+\textcolor{teal!80}{alice@dev:\mtilde}@+$ ./main
A number (0 to stop): 4
A number (0 to stop): -10
Negative numbers are ignored.
A number (0 to stop): 3
A number (0 to stop): 0
Total: 7
@+\textcolor{teal!80}{alice@dev:\mtilde}@+$ ./main
A number (0 to stop): 60
A number (0 to stop): 50
Total: 110
\end{lstlisting}

\begin{figure}[h]
  \centering
  \begin{adjustbox}{max width=\linewidth}
    \begin{tikzpicture}[node distance=6mm]

      \node[cfterm] (start) {start of main};
      \node[cfstep, below=of start] (init) {let mut total = 0;};
      \node[cfstep, below=8mm of init] (read) {let n = read!i32(...);};
      \node[cftest, below=of read] (zero) {n == 0};
      \node[cftest, below=of zero] (neg) {n < 0};
      \node[cfstep, below=of neg] (add) {total += n;};
      \node[cftest, below=of add] (full) {total >= 100};
      \node[cfstep, below=8mm of full] (print) {println("Total: ", total);};
      \node[cfterm, below=of print] (end) {end of main};
      \node[cfstep, right=12mm of neg] (ignored) {println("... ignored.");};

      \coordinate (out) at ($(full.west) + (-12mm, 0)$);
      \coordinate (back) at ($(ignored.east) + (8mm, 0)$);

      \draw[cfflow] (start) -- (init);
      \draw[cfflow] (init) -- node[cflab, right] {loop \{} (read);
      \draw[cfflow] (read) -- (zero);
      \draw[cfflow] (zero) -- node[cflab, right] {false} (neg);
      \draw[cfflow] (neg) -- node[cflab, right] {false} (add);
      \draw[cfflow] (add) -- (full);
      \draw[cfflow] (print) -- (end);

      \draw[cfflow] (zero.west) -- node[cflab, above] {true} (zero.west -| out)
        -- node[cfjump, left, pos=0.5] {break} (out) |- (print.west);
      \draw[cfflow] (full.west) -- node[cflab, above] {true} (out) |- (print.west);

      \draw[cfflow] (neg.east) -- node[cflab, above] {true} (ignored.west);
      \draw[cfflow] (ignored.north) -- node[cfjump, left] {continue} (ignored.north |- read.east)
        -- (read.east);
      \draw[cfflow] (full.east) -- node[cflab, above] {false} (full.east -| back)
        -- node[cflab, right, align=left] {end of\\the block} (back |- read.east)
        -- (read.east);
    \end{tikzpicture}
  \end{adjustbox}
  \caption{\label{fig:control:loop_jumps} The \token{loop} of
    Listing~\ref{lst:control:continue}. The end of the block and
    \token{continue} both go back to the top of the block; either
    \token{break} leaves the loop.}
\end{figure}


\token{break} and \token{continue} both make the lines that follow them in the
same block unreachable, since control never gets there. In the next listing,
the goodbye message on line~9 can never be printed: the \token{break} just
above it has already left the loop.

\begin{lstlisting}[style=coloredverbatimError, escapechar=@, caption=An instruction after a \token{break}]
use std::io;

fn main() {
  let mut total = 0;
  loop {
    let n = read!i32(ask-> "A number (0 to stop): ");
    if n == 0 {
      break;
      @\hb{println("Goodbye.")}@;
    }
    total += n;
  }
  println("Total: ", total);
}
\end{lstlisting}

The compiler rejects such lines with ``unreachable statement'', as it rejects a
condition that is always true: code that can never run is a mistake. Here the
message belongs before the \token{break}. For the same reason, both words are
errors outside of a loop, where there is nothing to leave and nothing to repeat.

\subsection{A loop that produces a value}

A \token{loop}, like an \token{if}, can be used as a value. The value comes from
the \token{break} that ends it: \token{break m;} leaves the loop, and \token{m}
becomes the value of the whole \token{loop}.

The next listing uses this to insist until the user gives a valid month. The
loop only ends through the \token{break} on line~7, which carries the month out,
so \token{month} is always between 1 and 12.

\begin{lstlisting}[style=coloredverbatim, caption=A loop whose value comes from \token{break}, label=lst:control:loop_value]
use std::io;

fn main() {
  let month = loop {
    let m = read!i32(ask-> "Month (1 to 12): ");
    if m >= 1 && m <= 12 {
      break m;
    }
    println("There is no month ", m, ".");
  };
  println("Month ", month, " it is.");
}
\end{lstlisting}
\vspace{-10pt}%
\begin{lstlisting}[style=bashVerb, escapechar=@+]
@+\textcolor{teal!80}{alice@dev:\mtilde}@+$ ./main
Month (1 to 12): 0
There is no month 0.
Month (1 to 12): 13
There is no month 13.
Month (1 to 12): 4
Month 4 it is.
\end{lstlisting}

The semicolon after the closing brace on line~10 ends the declaration of
\token{month}, as it would after any other value.

A \token{loop} with several \token{break}s takes its value from whichever one
ends it, so they must all carry a value, of the same type. The next listing
lets the user cancel by typing \token{0}, but its new \token{break} on line~7
carries nothing.

\begin{lstlisting}[style=coloredverbatimError, escapechar=@, caption=Two \token{break}s that do not agree]
use std::io;

fn main() {
  let month = loop {
    let m = read!i32(ask-> "Month (1 to 12, 0 to cancel): ");
    if m == 0 {
      break;
    }
    if m >= 1 && m <= 12 {
      @\hb{break m}@;
    }
    println("There is no month ", m, ".");
  };
  println("Month ", month, " it is.");
}
\end{lstlisting}

The compiler reports ``incompatible types void and i32'' on the second
\token{break}: the first one gave the loop no value, of type \token{void}, and
the second an \token{i32}. The cancellation needs a value of its own, such as
\token{break 0;}, which the rest of the program can then test.

\vfill\pagebreak
\section{Repeating while a condition holds: \texttt{while}}
\label{sec:control:while}

Many loops have the same shape: test a condition at the top, break if it does
not hold, otherwise do the work. \token{while} writes that shape directly.
\token{while condition \{ ... \}} tests \token{condition} before each iteration,
runs the block if it is \token{true}, and leaves the loop as soon as it is
\token{false}.

The next listing counts the digits of a number by dividing it by ten until a
single digit is left. Figure~\ref{fig:control:while} draws its flow: the arrow
from the bottom of the block goes back up to the test, not to the start of the
block.

\begin{lstlisting}[style=coloredverbatim, caption=Counting digits with a \token{while} loop, label=lst:control:while]
use std::io;

fn main() {
  let mut n = read!i32(ask-> "A positive number: ");
  let mut digits = 1;
  while n >= 10 {
    n = n / 10;
    digits += 1;
  }
  println("That number has ", digits, " digits.");
}
\end{lstlisting}
\vspace{-10pt}%
\begin{lstlisting}[style=bashVerb, escapechar=@+]
@+\textcolor{teal!80}{alice@dev:\mtilde}@+$ ./main
A positive number: 40953
That number has 5 digits.
\end{lstlisting}

\begin{figure}[h]
  \centering
  \begin{adjustbox}{max width=\linewidth}
    \begin{tikzpicture}[node distance=6mm]

      \node[cfterm] (start) {start of main};
      \node[cfstep, below=of start] (init) {let mut n = read!i32(...);\\let mut digits = 1;};
      \node[cftest, below=of init] (test) {n >= 10};
      \node[cfstep, below=8mm of test] (body) {n = n / 10;\\digits += 1;};
      \node[cfstep, below=of body] (print) {println(..., digits, ...);};
      \node[cfterm, below=of print] (end) {end of main};

      \draw[cfflow] (start) -- (init);
      \draw[cfflow] (init) -- (test);
      \draw[cfflow] (test) -- node[cflab, right] {true} (body);
      \draw[cfflow] (body.east) -- ++(14mm, 0) |- (test.east);
      \draw[cfflow] (test.west) -- node[cflab, above] {false} ++(-14mm, 0) |- (print.west);
      \draw[cfflow] (print) -- (end);
    \end{tikzpicture}
  \end{adjustbox}
  \caption{\label{fig:control:while} The \token{while} loop of
    Listing~\ref{lst:control:while}. The condition is tested before every
    iteration, including the first.}
\end{figure}


Because the test comes first, a \token{while} can run its block zero times: with
\token{7} as input, \token{n >= 10} is \token{false} from the start, and the
program reports a single digit without entering the loop. The same program written with
\token{loop} needs the test and the \token{break} spelled out:

\begin{lstlisting}[style=coloredverbatim, caption=The same loop without \token{while}]
use std::io;

fn main() {
  let mut n = read!i32(ask-> "A positive number: ");
  let mut digits = 1;
  loop {
    if !(n >= 10) {
      break;
    }
    n = n / 10;
    digits += 1;
  }
  println("That number has ", digits, " digits.");
}
\end{lstlisting}

\token{break} and \token{continue} work in a \token{while} exactly as in a
\token{loop}. A \token{while}, however, has no value: it may stop because its
condition became false, and then no \token{break} has provided one. So
\token{break} inside a \token{while} never carries a value.

\subsection{Choosing between \texttt{while} and \texttt{loop}}

A \token{while} whose condition is always \token{true} is a \token{loop} under
another name, and the compiler insists on the real name. \token{while true \{}
is rejected with ``infinite loop with always true test is not allowed, rather use
a loop construct''.

\begin{lstlisting}[style=coloredverbatimError, escapechar=@, caption=\token{while true} is written \token{loop}]
use std::io;

fn main() {
  let mut total = 0;
  while @\hb{true}@ {
    let n = read!i32(ask-> "A number (0 to stop): ");
    if n == 0 {
      break;
    }
    total += n;
  }
  println("Total: ", total);
}
\end{lstlisting}

Which of the two fits depends on where the loop decides to stop. When the
decision can be taken before the work, as in Listing~\ref{lst:control:while},
use \token{while}. When the work has to happen first -- the number must be read
before it can be compared with zero, the month must be asked before it can be
checked -- the decision sits in the middle or at the end of the block, and the
loop is a \token{loop} with a \token{break} there.

\notebox{Ymir has no \textit{do-while} loop, which some languages provide to test
  the condition after the block rather than before it. The \token{loop} covers
  that case, with the test written at the end of the block:
  \token{loop \{ ... if !condition \{ break; \} \}}.}

\vfill\pagebreak
\section{Counting: \texttt{for} and ranges}
\label{sec:control:for}

The third loop is for counting. \token{for i in 1 ... 5 \{ ... \}} runs its block
once for each integer from 1 to 5, in order, and inside the block the variable
\token{i} holds the current one. \token{i} is called the \textit{iterator} of
the loop, and \token{1 ... 5} is a \textit{range}.

\begin{lstlisting}[style=coloredverbatim, caption=Counting from 1 to 5, label=lst:control:for]
use std::io;

fn main() {
  for i in 1 ... 5 {
    println(i, " squared is ", i * i);
  }
}
\end{lstlisting}
\vspace{-10pt}%
\begin{lstlisting}[style=bashVerb]
1 squared is 1
2 squared is 4
3 squared is 9
4 squared is 16
5 squared is 25
\end{lstlisting}

\subsection{Ranges}

A range is written with two bounds and one of two operators between them.

\begin{itemize}
\item \token{a ... b}, with three dots, goes from \token{a} to \token{b}, both
  included: \token{1 ... 5} is 1, 2, 3, 4, 5.
\item \token{a .. b}, with two dots, stops just before \token{b}: \token{0 .. 5}
  is 0, 1, 2, 3, 4. It still counts five numbers, and a count from zero is so
  common that this is the form most loops use.
\end{itemize}

The bounds can be any integer expressions, including values read at run time.
They are computed once, before the first iteration. A range whose end comes
before its start, such as \token{10 .. 0}, is empty, and the loop runs zero
times; the next listing relies on that when the user types \token{0}.

\begin{lstlisting}[style=coloredverbatim, caption=A range whose end is read at run time, label=lst:control:for_sum]
use std::io;

fn main() {
  let n = read!i32(ask-> "Up to: ");
  let mut sum = 0;
  for i in 1 ... n {
    sum += i;
  }
  println("1 + 2 + ... + ", n, " = ", sum);
}
\end{lstlisting}
\vspace{-10pt}%
\begin{lstlisting}[style=bashVerb, escapechar=@+]
@+\textcolor{teal!80}{alice@dev:\mtilde}@+$ ./main
Up to: 100
1 + 2 + ... + 100 = 5050
\end{lstlisting}

To count down, call \token{.reverse()} on the range. \token{(1 ... 10).reverse()}
counts from 10 down to 1. The bound that \token{..} leaves out is still the
upper one, so \token{(0 .. 10).reverse()} counts from 9 down to 0, whereas
\token{(0 ... 10).reverse()} starts at 10. To count by steps larger than one,
call \token{.stepBy(k)}, which keeps one number out of every \token{k}. In both
cases the parentheses around the range are needed, so that the call applies to
the whole range rather than to its last bound.

\begin{lstlisting}[style=coloredverbatim, caption=Counting down and by steps, label=lst:control:for_reverse]
use std::io;

fn main() {
  for i in (1 ... 10).reverse() {
    print(i, " ");
  }
  println("liftoff!");

  for i in (0 ... 20).stepBy(5) {
    print(i, " ");
  }
  println();
}
\end{lstlisting}
\vspace{-10pt}%
\begin{lstlisting}[style=bashVerb]
10 9 8 7 6 5 4 3 2 1 liftoff!
0 5 10 15 20
\end{lstlisting}

\token{print}, used here for the first time, is \token{println} without the line
break at the end, so that the numbers stay on one line. A \token{println} with
no argument ends the line.

\subsection{The iterator}

The iterator is a constant. Each iteration gives it the next number of the
range, but the block cannot change it: \token{i = i + 1;} inside the loop is
rejected, because ``i is a value iterator, and cannot be used as a lvalue''. To
skip numbers, use \token{.stepBy} or \token{continue}.

The iterator is also a variable like any other, so the compiler rejects it when
the block never reads it. When only the number of repetitions matters, the
wildcard \token{\_} takes the place of its name. It declares nothing, and is
therefore never unused.

\begin{lstlisting}[style=coloredverbatim, caption=Repeating without reading the iterator, label=lst:control:for_wildcard]
use std::io;

fn main() {
  for _ in 0 .. 3 {
    println("Hip hip hooray!");
  }
}
\end{lstlisting}

\subsection{Loops inside loops}

The block of a loop can contain another loop. The inner loop then runs in full
during each iteration of the outer one. In the next listing the outer loop picks
a row and the inner loop prints the columns of that row.

A \token{break} leaves only the loop it is written in, the innermost one. On
line~7 it ends the current row as soon as the column goes past the row number,
and the outer loop carries on with the next row. That is what gives the output
its triangle shape. \token{continue} likewise affects only the innermost loop.

\begin{lstlisting}[style=coloredverbatim, caption=A \token{break} in a nested loop, label=lst:control:nested]
use std::io;

fn main() {
  for row in 1 ... 4 {
    for col in 1 ... 4 {
      if col > row {
        break;
      }
      print(row * col, " ");
    }
    println();
  }
}
\end{lstlisting}
\vspace{-10pt}%
\begin{lstlisting}[style=bashVerb]
1
2 4
3 6 9
4 8 12 16
\end{lstlisting}

\subsection{Which loop to use}

The three loops overlap, and choosing well makes a program easier to read.

\begin{center}
  \begin{adjustbox}{max width=\linewidth}
    \begin{tabular}{|l|l|}
      \hline
      When & Use\\[0pt]
      \hline
      \hline
      The number of iterations is known before the loop starts & \token{for}\\[0pt]
      The loop goes on as long as a condition, tested first, holds & \token{while}\\[0pt]
      The decision to stop comes after some of the work & \token{loop} and \token{break}\\[0pt]
      \hline
    \end{tabular}
  \end{adjustbox}
\end{center}

A later chapter shows that \token{for} does much more than count: it goes
through the elements of a collection in the same way it goes through the numbers
of a range.

\vfill\pagebreak
\section{A first look at \texttt{match}}
\label{sec:control:match}

A chain of \token{else if} that compares one value with several possibilities
reads badly: the value is repeated in every condition, and the eye has to check
that it really is the same value each time. \token{match} states the value once
and lists the possibilities below it.

\begin{lstlisting}[style=coloredverbatim, caption=Choosing on the value of a number, label=lst:control:match]
use std::io;

fn main() {
  let day = read!i32(ask-> "Day of the week (1 to 7): ");
  match day {
    6 | 7 => {
      println("Weekend.");
    }
    1 ... 5 => {
      println("Weekday.");
    }
    _ => {
      println("There is no day ", day, ".");
    }
  }
}
\end{lstlisting}
\vspace{-10pt}%
\begin{lstlisting}[style=bashVerb, escapechar=@+]
@+\textcolor{teal!80}{alice@dev:\mtilde}@+$ ./main
Day of the week (1 to 7): 6
Weekend.
@+\textcolor{teal!80}{alice@dev:\mtilde}@+$ ./main
Day of the week (1 to 7): 9
There is no day 9.
\end{lstlisting}

Each line of the form \token{pattern => \{ ... \}} is an \textit{arm}. The arms
are tried from top to bottom, as the conditions of an \token{else if} chain are,
and the first one whose pattern fits the value runs its block. The others are
skipped. The arms are not separated by commas.

The patterns of Listing~\ref{lst:control:match} show the three forms that work on
numbers:

\begin{itemize}
\item a value, such as \token{7}, which fits only that value; several values
  separated by \token{|} fit any one of them, so \token{6 | 7} is the weekend;
\item a range, such as \token{1 ... 5}, which fits every number the range
  contains, with the same two operators as in Section~\ref{sec:control:for};
\item the wildcard \token{\_}, which fits anything, and is the \token{else} of a
  \token{match}.
\end{itemize}

Since \token{\_} fits anything, an arm written after it could never run. The
compiler rejects such an arm with ``matcher expression is never evaluated''. The
catch-all arm always comes last.

\subsection{Choosing a value}

A \token{match}, like an \token{if}, is an expression, and each arm can produce
a value instead of running instructions. When an arm is a single value, its
braces can go.

\begin{lstlisting}[style=coloredverbatim, caption=A \token{match} used as a value, label=lst:control:match_value]
use std::io;

fn main() {
  let day = read!i32(ask-> "Day of the week (1 to 7): ");
  let name = match day {
    1 => "Monday"
    2 => "Tuesday"
    3 => "Wednesday"
    4 => "Thursday"
    5 => "Friday"
    6 => "Saturday"
    7 => "Sunday"
    _ => "no day at all"
  };
  println("Day ", day, " is ", name, ".");
}
\end{lstlisting}

The rule is the one of Section~\ref{sec:control:if}: every possible run must
leave \token{name} with a value. Without the last arm, a day of 9 would match
nothing, so a \token{match} used as a value must end with an arm that catches
everything. Leaving it out is an error: ``the pattern matching has no default
case''.

\subsection{Naming the value: bindings and guards}

An arm can give a name to the value it matches, so that its block can use it.
The name declares a variable, which holds the matched value and exists only in
that arm. It is written in one of two ways.

\begin{itemize}
\item \token{name = pattern} names the value when the pattern after \token{=}
  fits it. An arm whose pattern is \token{s = 1 ... 9} fits the numbers from 1
  to 9, and inside its block \token{s} is the one that matched.
\item A name alone is a pattern of its own. Like \token{\_}, it fits any value,
  but unlike \token{\_} it gives the arm's block that value to work with.
\end{itemize}

An arm can also test the value further. After the pattern, \token{if} followed
by a condition adds a \textit{guard}, and the arm fits only when the guard holds
too. The condition can use the name that the pattern declared.
Listing~\ref{lst:control:match_guard} uses a binding on line~10, and a name
with a guard on line~13.

\begin{lstlisting}[style=coloredverbatim, caption=Arms with a binding and a guard, label=lst:control:match_guard]
use std::io;

fn main() {
  let limit = read!i32(ask-> "Speed limit: ");
  let speed = read!i32(ask-> "Your speed: ");
  match speed {
    0 => {
      println("You are not moving.");
    }
    s = 1 ... 9 => {
      println("Only ", s, "? You could walk faster.");
    }
    s if s > limit => {
      println("Slow down, you are ", s - limit, " over the limit.");
    }
    _ => {
      println("All good.");
    }
  }
}
\end{lstlisting}
\vspace{-10pt}%
\begin{lstlisting}[style=bashVerb, escapechar=@+]
@+\textcolor{teal!80}{alice@dev:\mtilde}@+$ ./main
Speed limit: 50
Your speed: 62
Slow down, you are 12 over the limit.
@+\textcolor{teal!80}{alice@dev:\mtilde}@+$ ./main
Speed limit: 50
Your speed: 5
Only 5? You could walk faster.
\end{lstlisting}

A name in a pattern always declares a new variable; it never reads an existing
one, and that holds on either side of \token{=}. So an arm
\token{s = limit => \{ ... \}} does not mean ``when the speed equals the
limit''. It declares two variables, \token{s} and \token{limit}, both holding
the speed, and since the program already has a \token{limit}, the compiler stops
with ``declaration of limit shadows another declaration''. A comparison with a
variable goes in a guard: the arm that means ``equal to the limit'' is
\token{s if s == limit => \{ ... \}}.

This chapter matches integers only. \token{match} is far more capable than
that -- it can take values apart, as later chapters show -- and
Section~\ref{sec:flow:pattern_matching} lists everything it can do.

\vfill\pagebreak
\section{Putting it together: a guessing game}
\label{sec:control:guessing_game}

The program below picks a number between 1 and 100 and lets the user guess it,
answering each guess with ``too small'' or ``too big'' until the number is
found.

\begin{lstlisting}[style=coloredverbatim, caption=A guessing game, label=lst:control:guessing_game]
use std::io;
use std::rand;

fn main() {
  let secret = uniform(1, 100);
  let mut tries = 0;

  println("I am thinking of a number between 1 and 100.");
  loop {
    let guess = read!i32(ask-> "Your guess: ");
    tries += 1;

    if guess < secret {
      println("Too small.");
    } else if guess > secret {
      println("Too big.");
    } else {
      println("Found it in ", tries, " tries!");
      break;
    }
  }
}
\end{lstlisting}
\vspace{-10pt}%
\begin{lstlisting}[style=bashVerb, escapechar=@+]
@+\textcolor{teal!80}{alice@dev:\mtilde}@+$ gyc main.yr -o main
@+\textcolor{teal!80}{alice@dev:\mtilde}@+$ ./main
I am thinking of a number between 1 and 100.
Your guess: 50
Too small.
Your guess: 75
Too big.
Your guess: 62
Found it in 3 tries!
\end{lstlisting}

Read from the top, each construction has a job.

\begin{itemize}
\item \token{uniform} draws the secret on line~5, so it differs on every run.
  Every decision the program takes depends on it, or on what the user types.
\item \token{tries} is declared before the loop, so that it keeps counting from
  one guess to the next.
\item The \token{loop} on line~9 repeats the question. It is a \token{loop}
  rather than a \token{while}, because the program can only decide to stop after
  a guess has been read and compared.
\item The chain of \token{if}, \token{else if} and \token{else} sorts each
  guess into one of three cases. Exactly one of them runs.
\item The \token{break} on line~19, in the last case only, is the one way out
  of the loop.
\end{itemize}

\notebox{The session above guesses by halving: 50 is the middle of 1 to 100, 75
  the middle of what is left, and so on. Each answer halves the numbers still
  possible, and seven halvings take a hundred possibilities down to one, so a
  player who always halves never needs more than seven tries. The last exercise
  of the chapter builds on that.}

\vfill%
\pagebreak

\section*{Exercises}
\markright{\MakeUppercase{Exercises}}%
\subsection*{Exercise 1}

Write a program that prints the numbers from 1 to 15, one per line, except that
multiples of 3 are replaced by \textit{Fizz}, multiples of 5 by \textit{Buzz},
and multiples of both by \textit{FizzBuzz}.

\subsection*{Exercise 2}

What does the following program print? Work it out by hand, iteration by
iteration, then check your answer by compiling and running it.

\begin{lstlisting}[style=coloredverbatim]
use std::io;

fn main() {
  let mut i = 0;
  while i < 10 {
    i += 1;
    if i % 3 == 0 {
      continue;
    }
    if i == 8 {
      break;
    }
    print(i, " ");
  }
  println();
  println("i = ", i);
}
\end{lstlisting}

\subsection*{Exercise 3}

The following program finds the smallest number whose square is at least the
number typed by the user, and says whether the number typed is negative. It contains
three errors. Find them, correct them, and compile the program to test your
corrections.

\begin{lstlisting}[style=coloredverbatimError, caption=Ymir program with errors]
use std::io;

fn main() {
  let n = read!i32(ask-> "A number: ");
  let mut root = 0;
  while true {
    if root * root >= n {
      break;
    }
    root += 1;
  }

  let sign = if n < 0 { "negative" };
  if root {
    println("Smallest root: ", root);
  }
  println(n, " is ", sign);
}
\end{lstlisting}

\subsection*{Exercise 4}

Write a program that reads a positive number and prints the sum of its digits.
For \token{40953}, it prints \token{21}. (\textit{Hint: the remainder of a
division by ten is the last digit of a number, and the division by ten removes
that digit.})

\subsection*{Exercise 5}

Write a program that reads a height and draws a triangle of stars that many rows
high. For example, with a height of 4:

\begin{lstlisting}[style=bashVerb]
*
**
***
****
\end{lstlisting}

\subsection*{Exercise 6}

Write a program that reads a month number and prints how many days that month
has, ignoring leap years, or a message if there is no such month. Compute the
number of days with a single \token{match} used as a value. In what order must
the arms be written, and why?

\subsection*{Exercise 7}

Rewrite the second loop of Listing~\ref{lst:control:for_reverse}, which prints
\token{0 5 10 15 20}, with a \token{while} loop instead of a \token{for}.

\subsection*{Exercise 8}

Write a program that reads a number and tells whether it is prime, that is,
whether it is at least 2 and has no divisor other than 1 and itself.

\subsection*{Exercise 9}

Modify the guessing game of Listing~\ref{lst:control:guessing_game} so that the
player has only seven tries. After seven wrong guesses, the program stops and
reveals the number. (\textit{Hint: make the \token{loop} produce a \token{bool}
that says whether the number was found.})

\vfill%
\pagebreak

\forceevenpage{}

\section*{Solutions}
\markright{\MakeUppercase{Solutions}}%
\subsection*{Exercise 1}

The test for \textit{FizzBuzz} must come first: a multiple of 15 is also a
multiple of 3, and the chain stops at the first condition that holds.

\begin{lstlisting}[style=coloredverbatim, caption=Solution for exercise 1]
use std::io;

fn main() {
  for i in 1 ... 15 {
    if i % 15 == 0 {
      println("FizzBuzz");
    } else if i % 3 == 0 {
      println("Fizz");
    } else if i % 5 == 0 {
      println("Buzz");
    } else {
      println(i);
    }
  }
}
\end{lstlisting}

\subsection*{Exercise 2}

\begin{lstlisting}[style=bashVerb]
1 2 4 5 7
i = 8
\end{lstlisting}

3 and 6 are skipped by \token{continue}, which returns to the test of the
\token{while} before \token{print} is reached. When \token{i} reaches 8,
\token{break} leaves the loop, also before \token{print}, and \token{i} keeps the
value 8.

\subsection*{Exercise 3}

\token{while true} is not allowed, and is written \token{loop}. The \token{if}
that gives \token{sign} its value has no \token{else}, so there is no value when
\token{n} is not negative. And \token{root} is an integer, not a Boolean, so the
condition of the last \token{if} must be a comparison.

\begin{lstlisting}[style=coloredverbatim, escapechar=@, caption=Solution for exercise 3]
use std::io;

fn main() {
  let n = read!i32(ask-> "A number: ");
  let mut root = 0;
  @\hcb{loop}@ {
    if root * root >= n {
      break;
    }
    root += 1;
  }

  let sign = if n < 0 { "negative" } @\hcb{else \{ "positive or zero" \}}@;
  if @\hcb{root != 0}@ {
    println("Smallest root: ", root);
  }
  println(n, " is ", sign);
}
\end{lstlisting}

\subsection*{Exercise 4}

\begin{lstlisting}[style=coloredverbatim, caption=Solution for exercise 4]
use std::io;

fn main() {
  let mut n = read!i32(ask-> "A positive number: ");
  let mut sum = 0;
  while n > 0 {
    sum += n % 10;
    n /= 10;
  }
  println("Sum of the digits: ", sum);
}
\end{lstlisting}

\subsection*{Exercise 5}

The outer loop picks the row, the inner loop prints one star per column. The
inner iterator is never read, so it is \token{\_}.

\begin{lstlisting}[style=coloredverbatim, caption=Solution for exercise 5]
use std::io;

fn main() {
  let height = read!i32(ask-> "Height: ");
  for row in 1 ... height {
    for _ in 0 .. row {
      print("*");
    }
    println();
  }
}
\end{lstlisting}

\subsection*{Exercise 6}

The arms are tried from top to bottom. The arm \token{1 ... 12} also fits
February and the months of 30 days, so it must come after them, and the
catch-all \token{\_}, which fits everything, must come last. It is also required:
without it, the \token{match} would have no value for 13.

\begin{lstlisting}[style=coloredverbatim, caption=Solution for exercise 6]
use std::io;

fn main() {
  let month = read!i32(ask-> "Month (1 to 12): ");
  let days = match month {
    2 => 28
    4 | 6 | 9 | 11 => 30
    1 ... 12 => 31
    _ => 0
  };

  if days == 0 {
    println("There is no month ", month, ".");
  } else {
    println("Month ", month, " has ", days, " days.");
  }
}
\end{lstlisting}

\subsection*{Exercise 7}

The iterator becomes a mutable variable declared before the loop, the bound
becomes the condition, and the step becomes an addition at the end of the block.
The \token{for} did all three by itself, which is why it is the better choice
when the numbers are known in advance.

\begin{lstlisting}[style=coloredverbatim, caption=Solution for exercise 7]
use std::io;

fn main() {
  let mut i = 0;
  while i <= 20 {
    print(i, " ");
    i += 5;
  }
  println();
}
\end{lstlisting}

\subsection*{Exercise 8}

The program assumes the number is prime, then looks for a divisor between 2 and
the number itself, excluded. The first divisor found proves the assumption
wrong, and there is no need to look further. A number below 2 is not prime, and
the range \token{2 .. n} is then empty, so the loop does not run at all.

\begin{lstlisting}[style=coloredverbatim, caption=Solution for exercise 8]
use std::io;

fn main() {
  let n = read!i32(ask-> "A number: ");
  let mut prime = n >= 2;
  for d in 2 .. n {
    if n % d == 0 {
      prime = false;
      break;
    }
  }

  if prime {
    println(n, " is prime.");
  } else {
    println(n, " is not prime.");
  }
}
\end{lstlisting}

The loop tests more divisors than it needs to. A divisor of \token{n} other
than \token{n} itself is at most half of \token{n}, so the range can stop there,
\token{2 ... n / 2}, and the program does half the work for the same answer.

\subsection*{Exercise 9}

The loop has two ways out, and each \token{break} says which one was taken.
Checking the number of tries at the top of the block, before reading, stops the
game after the seventh wrong guess without asking for an eighth.

\begin{lstlisting}[style=coloredverbatim, caption=Solution for exercise 9]
use std::io;
use std::rand;

fn main() {
  let secret = uniform(1, 100);
  let mut tries = 0;

  println("I am thinking of a number between 1 and 100. You have 7 tries.");
  let found = loop {
    if tries == 7 {
      break false;
    }

    let guess = read!i32(ask-> "Your guess: ");
    tries += 1;

    if guess < secret {
      println("Too small.");
    } else if guess > secret {
      println("Too big.");
    } else {
      break true;
    }
  };

  if found {
    println("Found it in ", tries, " tries!");
  } else {
    println("Lost! The number was ", secret, ".");
  }
}
\end{lstlisting}

